\documentclass[11pt]{article}
\usepackage[margin=1in]{geometry}
\usepackage{amsmath}
\usepackage{graphicx}
\usepackage{hyperref}

% ==============================================================================
% STYLE GUIDELINES FOR THIS WHITEPAPER
% ==============================================================================
% Model: Bitcoin whitepaper (Satoshi Nakamoto, 2008)
%   - Length: ~9 pages total
%   - Tone: Understated confidence, no hype, matter-of-fact
%   - Structure: Problem → Solution → How it works → Math/Security → Conclusion
%   - Every sentence must earn its place
%
% FRAMING:
%   - Abstract: Technical (for blockchain researchers)
%   - Introduction: Society-level problems (for broader audience)
%     * Start with what people/businesses WANT but CAN'T HAVE due to trust
%     * Not "developers face constraints" but "businesses lose money to fraud"
%     * Bitcoin test: Would a VC/regulator/executive understand para 1?
%   - Technical sections: Precise but accessible
%
% TERMINOLOGY (use consistently):
%   - onchain / offchain (no hyphens)
%   - attestation (not "statement" or "proof" unless cryptographic)
%   - stake against / stake-backed (not "economically back")
%   - multi-step (not "temporal" or "multi-phase")
%   - validators (not "nodes" unless specifically P2P network context)
%   - fraud proofs (not "challenge mechanism")
%
% TONE RULES:
%   - "We propose..." not "We revolutionize..."
%   - State problems plainly: "X suffers from Y"
%   - No: "Imagine if...", "In the future...", "paradigm shift"
%   - Admit limitations honestly (like Bitcoin's privacy section)
%
% SATOSHI TEST:
%   - Would this sentence appear in the Bitcoin whitepaper?
%   - If too breathless/marketing-y, cut or rephrase
%   - If too verbose, tighten
%   - If academic filler ("various", "multiple", "different"), delete
% ==============================================================================

\title{Heterogeneous Verification for Multi-Step Decentralized Workflows}
\author{Cory Gabrielsen (cory.eth)}
\date{\today}

\begin{document}

\maketitle

\begin{abstract}
Financial and confidential automation relies on trusted intermediaries, imposing fees and custody risks that limit practical applications. Blockchains verify onchain operations but cannot verify automation requiring confidential execution, external data, or multi-step coordination. We propose a verification system where each operation uses the mechanism appropriate to its computation. The system combines cryptographic proofs for blockchain operations, hardware attestations for confidential execution, and stake-backed attestations for external interactions into workflow attestations. Validators broadcast attestations backed by stake, maintaining verifiability as long as honest validators control sufficient stake to make fraud economically irrational. This enables verifiable automation for financial operations, multi-party workflows, and confidential computation requiring trusted intermediaries.
\end{abstract}

\section{Introduction}

Automation of critical operations has come to rely almost exclusively on centralized platforms as trusted intermediaries to handle financial transactions, confidential data, and multi-party coordination. While these platforms work well enough for simple workflows, they suffer from the inherent vulnerabilities of the trust-based model. Automated processes that handle value cannot avoid custody risk, since platforms must hold private keys and can be hacked, bankrupted, or compromised. The cost of custody appears in billions of dollars lost to exchange failures, limiting the scale at which financial automation becomes viable and cutting off entire classes of processes requiring significant value custody. Beyond financial risk, confidential processes cannot be automated trustlessly—healthcare workflows expose patient data to cloud providers, enterprise systems surrender control to vendors, sensitive operations reveal private parameters to platforms. With platform dependencies comes concentration of power. Organizations must accept vendor lock-in and platform rules, granting more access than they would otherwise. A certain percentage of custody losses, data breaches, and censorship is accepted as unavoidable. These risks can be avoided through manual processes with trusted relationships, but no mechanism exists for trustless automation at scale without centralized operators.

What is needed is a verification system based on cryptographic attestations instead of trust, allowing automated workflows to execute across multiple steps without centralized intermediaries. Workflows combining cryptographic proofs for blockchain operations, hardware attestations for confidential execution, and stake-backed attestations for external interactions would protect against custody risk while preserving confidentiality and maintaining verifiability across heterogeneous trust requirements. In this paper, we propose a heterogeneous verification system where validators attest to each workflow step using mechanisms appropriate to the computation and stake economic value against their attestations. Fraudulent attestations can be challenged through onchain fraud proofs. The system is secure as long as honest validators control sufficient stake to make fraud economically irrational.

\section*{References}

\begin{enumerate}
\item[{[1]}] S. Nakamoto, ``Bitcoin: A Peer-to-Peer Electronic Cash System,'' 2008.
\end{enumerate}

\end{document}

